\addtotoc{Abstract}  % Add the "Abstract" page entry to the Contents
\abstract{
\addtocontents{toc}{\vspace{1em}}  % Add a gap in the Contents, for aesthetics

The rapid proliferation of multimodal misinformation, where images are paired with misleading text to fabricate credibility, poses a growing challenge that exceeds the capacity of manual fact-checking. This thesis presents a hybrid multimodal retrieval-augmented pipeline for automated fact-checking of image-text claims. The system uses sparse lexical retrieval (SPLADE), dense semantic retrieval (BGE), and long-context visual retrieval (LongCLIP) with post-retrieval refinement through cross-encoder reranking and Maximal Marginal Relevance (MMR) diversification. A Vision-Language Model (Qwen3-VL-8B-Instruct) generates verification questions, extracts answers from retrieved evidence, and produces grounded verdicts with justifications.

We evaluate the pipeline on the AVerImaTeC shared task, a multimodal fact verification benchmark containing real-world claims from professional fact-checking websites. Through incremental ablation, we demonstrate that LongCLIP improves evidence retrieval by 29\% over standard CLIP and that adding MMR diversification provides the largest scores across all metrics. A comparison of question generation strategies shows that claim based retrieval for Question-Answer generation substantially outperforms iterative approaches by avoiding error propagation. We further reveal that in-context learning exposes a generalization gap: it reduces validation accuracy on the heavily imbalanced dataset but improves generalization to unseen test data by mitigating majority-class bias.

On the validation set, our system achieves the highest verdict score among all participating teams. On the test set, it outperforms the official baseline across all metrics and achieves competitive performance relative to systems using proprietary models orders of magnitude larger, while running entirely on a single open-source 8B-parameter model on one GPU. These results demonstrate that well-designed retrieval system can substantially compensate for smaller reasoning models in multimodal fact-checking.

}
